% Auto-generated from raw.txt
% Usage:
%   \vocabentry{word}{/ipa/ (pos)}{definition}{example}{note}

\vocabentry{Cognitive}
{/ˈkɑːɡnətɪv/ (adj)}
{Relating to the mental action or process of acquiring knowledge and understanding.}
{Cognitive development in children is influenced by both genetics and environment.}
{Đây là một tính từ dùng để chỉ các quá trình thuộc về nhận thức như tư duy, trí nhớ và lý luận.}

\vocabentry{Ambivert}
{/ˈæmbivɜːrt/ (n)}
{A person whose personality has a balance of extrovert and introvert features.}
{As an ambivert, she enjoys social gatherings but also values her alone time to recharge.}
{Đây là một danh từ dùng để chỉ người có tính cách hướng trung, pha trộn giữa hướng nội và hướng ngoại.}

\vocabentry{Conscientious}
{/ˌkɑːnʃiˈenʃəs/ (adj)}
{Wishing to do one's work or duty well and thoroughly.}
{A conscientious student always meets deadlines and pays attention to the smallest details.}
{Đây là một tính từ dùng để chỉ người tận tâm, chu đáo và luôn làm việc có trách nhiệm.}

\vocabentry{Resilience}
{/rɪˈzɪliəns/ (n)}
{The capacity to recover quickly from difficulties; toughness.}
{Psychological resilience helps individuals cope with stress and trauma effectively.}
{Đây là một danh từ dùng để chỉ khả năng phục hồi, sự kiên cường trước nghịch cảnh.}

\vocabentry{Altruism}
{/ˈæltruɪzəm/ (n)}
{The belief in or practice of disinterested and selfless concern for the well-being of others.}
{Many charities rely on the altruism of donors to fund their community projects.}
{Đây là một danh từ dùng để chỉ lòng vị tha, sự quan tâm không vụ lợi đến người khác.}

\vocabentry{Neuroticism}
{/njʊˈrɒtɪsɪzəm/ (n)}
{A personality trait characterized by anxiety, fear, moodiness, and physical complaints.}
{High levels of neuroticism can lead to frequent feelings of worry and emotional instability.}
{Đây là một danh từ dùng để chỉ tâm lý bất ổn, dễ bị lo âu hoặc căng thẳng.}

\vocabentry{Empathy}
{/ˈempəθi/ (n)}
{The ability to understand and share the feelings of another.}
{Doctors need a high degree of empathy to provide better care for their patients.}
{Đây là một danh từ dùng để chỉ sự đồng cảm, khả năng thấu hiểu cảm xúc của người khác.}

\vocabentry{Narcissism}
{/ˈnɑːrsɪsɪzəm/ (n)}
{Excessive interest in or admiration of oneself and one's physical appearance.}
{Social media is often criticized for fueling narcissism among young adults.}
{Đây là một danh từ dùng để chỉ ái kỷ, sự tự phụ hoặc quá yêu bản thân.}

\vocabentry{Extroversion}
{/ˌekstrəˈvɜːrʒn/ (n)}
{The quality of being outgoing and directing attention to things other than oneself.}
{His extroversion makes him an excellent candidate for sales and public relations roles.}
{Đây là một danh từ dùng để chỉ tính hướng ngoại.}

\vocabentry{Introversion}
{/ˌɪntrəˈvɜːrʒn/ (n)}
{The state of or tendency toward being wholly or predominantly concerned with and interested in one's own mental life.}
{Introversion should not be confused with shyness, as it is simply a preference for quiet environments.}
{Đây là một danh từ dùng để chỉ tính hướng nội.}

\vocabentry{Subconscious}
{/ˌsʌbˈkɑːnʃəs/ (adj/n)}
{Concerning the part of the mind of which one is not fully aware but which influences one's actions.}
{Many of our fears are rooted in our subconscious mind and childhood experiences.}
{Đây là một tính từ/danh từ dùng để chỉ tiềm thức.}

\vocabentry{Behaviorism}
{/bɪˈheɪvjərɪzəm/ (n)}
{The theory that human behavior can be explained in terms of conditioning.}
{Behaviorism focuses on observable actions rather than internal mental states.}
{Đây là một danh từ dùng để chỉ thuyết hành vi trong tâm lý học.}

\vocabentry{Temperament}
{/ˈtemprəmənt/ (n)}
{A person's or animal's nature, especially as it permanently affects their behavior.}
{My sister has a very calm temperament, even when dealing with stressful situations.}
{Đây là một danh từ dùng để chỉ khí chất hoặc tính khí tự nhiên của một người.}

\vocabentry{Phobia}
{/ˈfoʊbiə/ (n)}
{An extreme or irrational fear of or aversion to something.}
{Claustrophobia is a common phobia involving the fear of being in enclosed spaces.}
{Đây là một danh từ dùng để chỉ nỗi ám ảnh cưỡng chế hoặc nỗi sợ vô lý.}

\vocabentry{Perception}
{/pərˈsepʃn/ (n)}
{The ability to see, hear, or become aware of something through the senses.}
{Public perception of mental health issues has changed significantly over the last decade.}
{Đây là một danh từ dùng để chỉ sự nhận thức hoặc cách nhìn nhận sự việc.}

\vocabentry{Self-esteem}
{/ˌself ɪˈstiːm/ (n)}
{Confidence in one's own worth or abilities; self-respect.}
{Low self-esteem can prevent people from pursuing their goals and dreams.}
{Đây là một danh từ dùng để chỉ lòng tự trọng hoặc sự tự tin vào giá trị bản thân.}

\vocabentry{Motivation}
{/ˌmoʊtɪˈveɪʃn/ (n)}
{The reason or reasons one has for acting or behaving in a particular way.}
{Intrinsic motivation is often more powerful than external rewards for long-term success.}
{Đây là một danh từ dùng để chỉ động lực.}

\vocabentry{Trauma}
{/ˈtraʊmə/ (n)}
{A deeply distressing or disturbing experience.}
{It takes time and professional support to recover from childhood trauma.}
{Đây là một danh từ dùng để chỉ chấn thương tâm lý hoặc cú sốc tinh thần.}

\vocabentry{Anxiety}
{/æŋˈzaɪəti/ (n)}
{A feeling of worry, nervousness, or unease, typically about an imminent event.}
{Test anxiety can prevent students from performing to the best of their abilities.}
{Đây là một danh từ dùng để chỉ sự lo âu hoặc trạng thái bồn chồn.}

\vocabentry{Cognitive dissonance}
{/ˌkɑːɡnətɪv ˈdɪsənəns/ (n)}
{The state of having inconsistent thoughts, beliefs, or attitudes.}
{People often experience cognitive dissonance when their actions do not align with their values.}
{Đây là một cụm danh từ dùng để chỉ sự mâu thuẫn nhận thức.}

\vocabentry{Archetype}
{/ˈɑːrkitaɪp/ (n)}
{A very typical example of a certain person or thing.}
{The "Hero" is a common archetype found in the myths and legends of almost every culture.}
{Trong tâm lý học, đây là danh từ chỉ các hình mẫu nguyên mẫu trong tâm trí con người.}

\vocabentry{Conditioning}
{/kənˈdɪʃənɪŋ/ (n)}
{The process of training or accustoming a person or animal to behave in a certain way.}
{Pavlov's experiments are the most famous examples of classical conditioning.}
{Đây là một danh từ dùng để chỉ sự điều kiện hóa/phản xạ có điều kiện.}

\vocabentry{Defense mechanism}
{/dɪˈfens ˈmekənɪzəm/ (n)}
{An automatic mental process used to avoid conscious conflict or anxiety.}
{Denial is a common defense mechanism used to protect oneself from painful truths.}
{Đây là một cụm danh từ dùng để chỉ cơ chế phòng vệ tâm lý.}

\vocabentry{Ego}
{/ˈiːɡoʊ/ (n)}
{A person's sense of self-esteem or self-importance.}
{He has a massive ego and finds it very difficult to accept any kind of criticism.}
{Đây là một danh từ dùng để chỉ cái tôi.}

\vocabentry{Empathize}
{/ˈempəθaɪz/ (v)}
{Understand and share the feelings of another.}
{It is easier to empathize with someone if you have experienced a similar situation.}
{Đây là một động từ dùng để chỉ việc đồng cảm với ai đó.}

\vocabentry{Gratification}
{/ˌɡrætɪfɪˈkeɪʃn/ (n)}
{Pleasure, especially when gained from the satisfaction of a desire.}
{The ability to delay gratification is often a predictor of future success.}
{Đây là một danh từ dùng để chỉ sự hài lòng hoặc thỏa mãn.}

\vocabentry{Hallucination}
{/həˌluːsɪˈneɪʃn/ (n)}
{An experience involving the apparent perception of something not present.}
{Certain mental illnesses can cause patients to suffer from vivid visual hallucinations.}
{Đây là một danh từ dùng để chỉ ảo giác.}

\vocabentry{Inhibition}
{/ˌɪnhɪˈbɪʃn/ (n)}
{A feeling that makes one self-conscious and unable to act in a relaxed and natural way.}
{Alcohol often lowers people's inhibitions, making them more talkative and outgoing.}
{Đây là một danh từ dùng để chỉ sự ức chế hoặc sự e dè, thiếu tự nhiên.}

\vocabentry{Introvert}
{/ˈɪntrəvɜːrt/ (n)}
{A shy, reticent person.}
{As an introvert, he prefers small gatherings of close friends over large parties.}
{Đây là một danh từ dùng để chỉ người hướng nội.}

\vocabentry{Libido}
{/lɪˈbiːdoʊ/ (n)}
{Sexual desire.}
{Freud believed that the libido was a primary driving force behind human behavior.}
{Trong tâm lý học, đây là danh từ chỉ năng lượng tình dục hoặc bản năng sống.}

\vocabentry{Masochism}
{/ˈmæsəkɪzəm/ (n)}
{The tendency to derive pleasure from one's own pain or humiliation.}
{His constant pursuit of difficult and painful tasks seemed almost like a form of masochism.}
{Đây là một danh từ dùng để chỉ chứng bạo dâm hoặc xu hướng tự ngược đãi.}

\vocabentry{Melancholy}
{/ˈmelənkɑːli/ (n/adj)}
{A feeling of pensive sadness, typically with no obvious cause.}
{The rainy weather always puts me in a state of deep melancholy.}
{Đây là một danh từ/tính từ dùng để chỉ sự u sầu hoặc trầm mặc.}

\vocabentry{Obsession}
{/əbˈseʃn/ (n)}
{A persistent, often uncontrollable, preoccupation with an idea, feeling, or person.}
{His obsession with detail made him a very effective, but often stressed, accountant.}
{Đây là một danh từ dùng để chỉ sự ám ảnh/nỗi ám ảnh.}

\vocabentry{Paranoia}
{/ˌpærəˈnɔɪə/ (n)}
{A mental condition characterized by delusions of persecution.}
{Suffering from severe paranoia, he was convinced that his neighbors were spying on him.}
{Đây là một danh từ dùng để chỉ chứng hoang tưởng bộ phát.}

\vocabentry{Projection}
{/prəˈdʒekʃn/ (n)}
{The mental process by which people attribute to others what is in their own minds.}
{Her anger towards her boss was actually a projection of her own feelings of inadequacy.}
{Trong tâm lý học, đây là sự phóng chiếu cảm xúc cá nhân lên người khác.}

\vocabentry{Psychosis}
{/saɪˈkoʊsɪs/ (n)}
{A severe mental disorder in which thought and emotions are so impaired that contact is lost with external reality.}
{During a period of psychosis, the patient was unable to distinguish between reality and fantasy.}
{Đây là một danh từ dùng để chỉ bệnh tâm thần phân liệt hoặc rối loạn tâm thần nặng.}

\vocabentry{Rationalization}
{/ˌræʃnələˈzeɪʃn/ (n)}
{The action of attempting to explain or justify behavior or an attitude with logical reasons.}
{His excuse for failing the test was just a rationalization of his lack of study.}
{Đây là một danh từ dùng để chỉ sự hợp lý hóa, tìm cớ bào chữa cho hành vi.}

\vocabentry{Regression}
{/rɪˈɡreʃn/ (n)}
{A return to a former or less developed state.}
{The child showed signs of regression, such as thumb-sucking, after the birth of the new baby.}
{Trong tâm lý học, đây là sự thoái lui về hành vi của lứa tuổi nhỏ hơn khi gặp căng thẳng.}

\vocabentry{Repression}
{/rɪˈpreʃn/ (n)}
{The action or process of suppressing a thought or desire in oneself so that it remains unconscious.}
{The repression of painful memories can sometimes lead to psychological problems later in life.}
{Đây là một danh từ dùng để chỉ sự dồn nén cảm xúc vào vô thức.}

\vocabentry{Sadism}
{/ˈseɪdɪzəm/ (n)}
{The tendency to derive pleasure from inflicting pain, suffering, or humiliation on others.}
{The villain in the movie showed a cruel sense of sadism towards his captives.}
{Đây là một danh từ dùng để chỉ chứng ác dâm hoặc thú vui làm khổ người khác.}

\vocabentry{Schizophrenia}
{/ˌskɪtsəˈfriːniə/ (n)}
{A long-term mental disorder involving a breakdown in the relation between thought, emotion, and behavior.}
{Schizophrenia is a complex condition that often requires long-term medication and support.}
{Đây là một danh từ dùng để chỉ bệnh tâm thần phân liệt.}

\vocabentry{Self-actualization}
{/ˌself ˌæktʃuələˈzeɪʃn/ (n)}
{The realization or fulfillment of one's talents and potentialities.}
{Maslow believed that self-actualization is the ultimate goal of human psychological development.}
{Đây là một danh từ dùng để chỉ sự tự hiện thực hóa, đỉnh cao trong tháp nhu cầu Maslow.}

\vocabentry{Socialization}
{/ˌsoʊʃələˈzeɪʃn/ (n)}
{The process of learning to behave in a way that is acceptable to society.}
{Schools play a vital role in the socialization of young children.}
{Đây là một danh từ dùng để chỉ quá trình xã hội hóa.}

\vocabentry{Sublimation}
{/ˌsʌblɪˈmeɪʃn/ (n)}
{The diversion of the energy of a biological or instinctive impulse into a socially acceptable activity.}
{He used exercise as a form of sublimation for his aggressive impulses.}
{Đây là một danh từ dùng để chỉ sự thăng hoa cảm xúc, chuyển đổi năng lượng tiêu cực thành hành động tích cực.}

\vocabentry{Suppression}
{/səˈpreʃn/ (n)}
{The action of suppressing something such as an activity or publication.}
{Conscious suppression of anger can sometimes lead to physical health issues like high blood pressure.}
{Trong tâm lý, đây là sự kìm nén cảm xúc một cách có ý thức.}

\vocabentry{Transference}
{/ˈtrænsfərəns/ (n)}
{The redirection of feelings and desires, especially of those unconsciously retained from childhood, toward a new object.}
{The therapist noticed that the patient was exhibiting signs of transference.}
{Đây là một danh từ dùng để chỉ sự chuyển di cảm xúc, thường xảy ra giữa bệnh nhân và bác sĩ.}

\vocabentry{Unconscious}
{/ʌnˈkɑːnʃəs/ (adj/n)}
{The part of the mind which is inaccessible to the conscious mind but which affects behavior and emotions.}
{Freud explored the role of the unconscious mind in shaping human personality.}
{Đây là một tính từ/danh từ dùng để chỉ vô thức.}

\vocabentry{Withdrawal}
{/wɪðˈdrɔːəl/ (n)}
{The action of withdrawing something; or social isolation.}
{Social withdrawal can be an early warning sign of depression.}
{Trong tâm lý, đây là sự thu mình hoặc rút lui khỏi các mối quan hệ xã hội.}

\vocabentry{Affective}
{/əˈfektɪv/ (adj)}
{Relating to moods, feelings, and attitudes.}
{Bipolar disorder is classified as an affective disorder.}
{Đây là một tính từ dùng để chỉ các yếu tố thuộc về cảm xúc hoặc tâm trạng.}

\vocabentry{Affiliation}
{/əˌfɪliˈeɪʃn/ (n)}
{The state or process of affiliating or being affiliated.}
{Humans have a strong natural need for social affiliation and belonging.}
{Trong tâm lý là nhu cầu liên kết hoặc thuộc về một nhóm nào đó.}

\vocabentry{Aggression}
{/əˈɡreʃn/ (n)}
{Feelings of anger or antipathy resulting in hostile or violent behavior.}
{Video games are often debated for their potential role in increasing aggression in children.}
{Đây là một danh từ dùng để chỉ sự hung tính hoặc hành vi xâm lược.}

\vocabentry{Amnesia}
{/æmˈniːʒə/ (n)}
{A partial or total loss of memory.}
{After the car accident, he suffered from temporary amnesia and couldn't remember his own name.}
{Đây là một danh từ dùng để chỉ chứng mất trí nhớ.}

\vocabentry{Analysis}
{/əˈnæləsɪs/ (n)}
{Detailed examination of the elements or structure of something.}
{Psychoanalysis is a method of therapy that involves the deep analysis of dreams and memories.}
{Trong tâm lý học là sự phân tâm hoặc phân tích tâm lý.}

\vocabentry{Apathy}
{/ˈæpəθi/ (n)}
{Lack of interest, enthusiasm, or concern.}
{Voter apathy is a major problem in many modern democracies.}
{Đây là một danh từ dùng để chỉ sự thờ ơ, vô cảm.}

\vocabentry{Assertiveness}
{/əˈsɜːrtɪvnəs/ (n)}
{Confident and forceful behavior.}
{Training in assertiveness can help people communicate their needs more effectively.}
{Đây là một danh từ dùng để chỉ sự quyết đoán hoặc khả năng tự khẳng định mình.}

\vocabentry{Attachment}
{/əˈtætʃmənt/ (n)}
{An extra part or extension that is or can be attached to something.}
{Early childhood attachment styles can influence future romantic relationships.}
{Trong tâm lý học là sự gắn kết tình cảm, đặc biệt giữa trẻ và mẹ.}

\vocabentry{Attitude}
{/ˈætɪtuːd/ (n)}
{A settled way of thinking or feeling about someone or something.}
{Having a positive attitude can make a big difference when facing difficult challenges.}
{Đây là một danh từ dùng để chỉ thái độ.}

\vocabentry{Avoidance}
{/əˈvɔɪdəns/ (n)}
{The action of keeping away from or not doing something.}
{Avoidance behavior is common among people suffering from social anxiety.}
{Trong tâm lý là hành vi né tránh những điều gây khó chịu.}

\vocabentry{Catharsis}
{/kəˈθɑːrsɪs/ (n)}
{The process of releasing, and thereby providing relief from, strong or repressed emotions.}
{Writing in a journal can be a powerful form of catharsis for many people.}
{Đây là một danh từ dùng để chỉ sự giải tỏa cảm xúc, sự thanh tẩy tâm hồn.}

\vocabentry{Clairvoyance}
{/klerˈvɔɪəns/ (n)}
{The supposed faculty of perceiving things or events in the future or beyond normal sensory contact.}
{She claimed to have the gift of clairvoyance and could predict future events.}
{Đây là một danh từ dùng để chỉ khả năng thấu thị hoặc ngoại cảm.}

\vocabentry{Compulsion}
{/kəmˈpʌlʃn/ (n)}
{An irresistible urge to behave in a certain way.}
{He felt a strong compulsion to check that the door was locked every five minutes.}
{Đây là một danh từ dùng để chỉ hành vi cưỡng chế, thôi thúc không cưỡng lại được.}

\vocabentry{Conflict}
{/ˈkɑːnflɪkt/ (n)}
{A serious disagreement or argument.}
{Inner conflict between duty and desire can cause a great deal of stress.}
{Trong tâm lý học là sự xung đột nội tâm.}

\vocabentry{Conformity}
{/kənˈfɔːrməti/ (n)}
{Compliance with standards, rules, or laws.}
{Social conformity often pressures individuals to act like everyone else.}
{Đây là một danh từ dùng để chỉ sự tuân thủ hoặc tính đồng nhất với đám đông.}

\vocabentry{Consensus}
{/kənˈsensəs/ (n)}
{A general agreement.}
{The team reached a consensus after a long and difficult discussion.}
{Sự đồng thuận trong nhóm.}

\vocabentry{Correlation}
{/ˌkɔːrəˈleɪʃn/ (n)}
{A mutual relationship or connection between two or more things.}
{There is a strong correlation between poverty and poor mental health.}
{Đây là một danh từ dùng để chỉ sự tương quan giữa các yếu tố tâm lý.}

\vocabentry{Counseling}
{/ˈkaʊnsəlɪŋ/ (n)}
{The provision of assistance and guidance in resolving personal or psychological problems.}
{Marriage counseling can help couples improve their communication and resolve conflicts.}
{Đây là một danh từ dùng để chỉ hoạt động tư vấn tâm lý.}

\vocabentry{Creativity}
{/ˌkriːeɪˈtɪvəti/ (n)}
{The use of the imagination or original ideas, especially in the production of an artistic work.}
{Creativity is a highly valued trait in many different fields of work.}
{Đây là một danh từ dùng để chỉ sự sáng tạo.}

\vocabentry{Crisis}
{/ˈkraɪsɪs/ (n)}
{A time of intense difficulty, trouble, or danger.}
{A midlife crisis often involves a period of intense self-reflection and change.}
{Đây là một danh từ dùng để chỉ cuộc khủng hoảng tâm lý.}

\vocabentry{Critical}
{/ˈkrɪtɪkl/ (adj)}
{Expressing adverse or disapproving comments or judgments.}
{Being overly critical of yourself can lead to low self-esteem.}
{Đây là một tính từ dùng để chỉ sự chỉ trích hoặc mang tính phê phán.}

\vocabentry{Delusion}
{/dɪˈluːʒn/ (n)}
{An idiosyncratic belief or impression that is firmly maintained despite being contradicted by what is generally accepted as reality.}
{He suffered from the delusion that he was a famous movie star.}
{Đây là một danh từ dùng để chỉ sự hoang tưởng.}

\vocabentry{Dependence}
{/rɪˈlaɪəns/ (n)}
{The state of relying on or being controlled by someone or something else.}
{Emotional dependence on a partner can be unhealthy for both individuals.}
{Sự lệ thuộc tâm lý.}

\vocabentry{Depression}
{/dɪˈpreʃn/ (n)}
{A mental health disorder characterized by persistently depressed mood or loss of interest in activities.}
{Depression is more than just feeling sad; it is a serious medical condition.}
{Đây là một danh từ dùng để chỉ bệnh trầm cảm.}

\vocabentry{Desire}
{/dɪˈzaɪər/ (n/v)}
{A strong feeling of wanting to have something or wishing for something to happen.}
{The desire for social approval can influence many of our daily decisions.}
{Đây là một danh từ/động từ dùng để chỉ khao khát hoặc ham muốn.}

\vocabentry{Development}
{/dɪˈveləpmənt/ (n)}
{The process of developing or being developed.}
{Psychological development continues throughout a person's entire life.}
{Sự phát triển tâm lý.}

\vocabentry{Deviation}
{/ˌdiːviˈeɪʃn/ (n)}
{The action of departing from an established course or accepted standard.}
{Any significant deviation from normal behavior should be carefully investigated.}
{Sự lệch lạc về hành vi hoặc tâm lý.}

\vocabentry{Disability}
{/ˌdɪsəˈbɪləti/ (n)}
{A physical or mental condition that limits a person's movements, senses, or activities.}
{Learning disabilities like dyslexia can make school very challenging for some children.}
{Đây là một danh từ dùng để chỉ sự khiếm khuyết hoặc khuyết tật.}

\vocabentry{Discrimination}
{/dɪˌskrɪmɪˈneɪʃn/ (n)}
{The unjust or prejudicial treatment of different categories of people or things.}
{Racial discrimination can have a devastating impact on a person's mental health.}
{Đây là một danh từ dùng để chỉ sự phân biệt đối xử.}

\vocabentry{Displacement}
{/dɪsˈpleɪsmənt/ (n)}
{The moving of something from its place or position.}
{Kicking the dog after a bad day at work is a classic example of displacement.}
{Trong tâm lý là sự dời chỗ cảm xúc, trút giận lên đối tượng khác.}

\vocabentry{Dissonance}
{/ˈdɪsənəns/ (n)}
{Lack of harmony or agreement.}
{Dissonance between a person's beliefs and actions can cause psychological discomfort.}
{Sự không hòa hợp, thường dùng trong "cognitive dissonance".}

\vocabentry{Distortion}
{/dɪˈstɔːrʃn/ (n)}
{The action of giving a misleading account or impression.}
{Cognitive distortion can lead people to view the world in a very negative way.}
{Sự bóp méo nhận thức hoặc sự thật.}

\vocabentry{Dream}
{/driːm/ (n/v)}
{A series of thoughts, images, and sensations occurring in a person's mind during sleep.}
{Freud believed that dreams were the "royal road to the unconscious."}
{Đây là một danh từ/động từ chỉ giấc mơ.}

\vocabentry{Drive}
{/draɪv/ (n/v)}
{An innate biologically determined urge to attain a goal or satisfy a need.}
{The drive for survival is the most basic human instinct.}
{Đây là một danh từ dùng để chỉ bản năng thôi thúc hoặc động lực.}

\vocabentry{Dynamic}
{/daɪˈnæmɪk/ (adj)}
{(Of a process or system) characterized by constant change, activity, or progress.}
{Family dynamics can be very complex and influence a child's personality.}
{Đây là một tính từ dùng để chỉ tính năng động.}

\vocabentry{Dysfunction}
{/dɪsˈfʌŋkʃn/ (n)}
{Abnormality or impairment in the function of a specified bodily organ or system.}
{Family dysfunction can lead to long-term emotional problems for children.}
{Sự rối loạn chức năng tâm lý.}

\vocabentry{Emotion}
{/ɪˈmoʊʃn/ (n)}
{A strong feeling deriving from one's circumstances, mood, or relationships with others.}
{Managing your emotions is a key part of emotional intelligence.}
{Đây là một danh từ dùng để chỉ cảm xúc.}

\vocabentry{Environment}
{/ɪnˈvaɪrənmənt/ (n)}
{The surroundings or conditions in which a person, animal, or plant lives or operates.}
{A supportive home environment is essential for a child's healthy development.}
{Môi trường sống ảnh hưởng đến tâm lý.}

\vocabentry{Equilibrium}
{/ˌiːkwɪˈlɪbriəm/ (n)}
{A state in which opposing forces or influences are balanced.}
{Meditation can help people regain their emotional equilibrium after a stressful event.}
{Trạng thái cân bằng tâm trí.}

\vocabentry{Eros}
{/ˈerɑːs/ (n)}
{The god of love; or life instinct in Freudian theory.}
{Eros represents the creative and life-preserving instincts in human psychology.}
{Bản năng sống/yêu thương.}

\vocabentry{Ethics}
{/ˈeθɪks/ (n)}
{Moral principles that govern a person's behavior or the conducting of an activity.}
{Psychologists must follow strict professional ethics when working with clients.}
{Đây là một danh từ dùng để chỉ đạo đức.}

\vocabentry{Evolution}
{/ˌevəˈluːʃn/ (n)}
{The gradual development of something, especially from a simple to a more complex form.}
{Evolutionary psychology explores how our ancient ancestors' behaviors still influence us today.}
{Đây là một danh từ dùng để chỉ sự tiến hóa, bao gồm cả tâm lý học tiến hóa.}

\vocabentry{Extinction}
{/ɪkˈstɪŋkʃn/ (n)}
{The state or process of a species, family, or larger group being or becoming extinct.}
{In behavioral therapy, extinction involves stopping a behavior by removing the reward.}
{Trong tâm lý học hành vi, đây là sự dập tắt một phản xạ.}

\vocabentry{Extravert}
{/ˈekstrəvɜːrt/ (n)}
{An outgoing, overtly expressive person.}
{Extraverts usually feel energized after spending time in large groups of people.}
{Đây là một danh từ dùng để chỉ người hướng ngoại.}

\vocabentry{Facilitation}
{/fəˌsɪlɪˈteɪʃn/ (n)}
{The action of making an action or process easy or easier.}
{Social facilitation occurs when people perform better on simple tasks when others are watching.}
{Sự tạo điều kiện thuận lợi.}

\vocabentry{Fantasy}
{/ˈfæntəsi/ (n)}
{The faculty or activity of imagining things, especially things that are impossible or improbable.}
{Many children spend a lot of time lost in a world of fantasy and imagination.}
{Đây là một danh từ dùng để chỉ sự tưởng tượng hoặc mộng tưởng.}

\vocabentry{Fear}
{/fɪr/ (n/v)}
{An unpleasant emotion caused by the belief that someone or something is dangerous.}
{Fear is a natural response that helps to protect us from danger.}
{Đây là một danh từ/động từ dùng để chỉ nỗi sợ hãi.}

\vocabentry{Fixation}
{/fɪkˈseɪʃn/ (n)}
{An obsessive interest in or feeling about someone or something.}
{His fixation with his physical appearance makes him very self-conscious.}
{Đây là một danh từ dùng để chỉ sự ám ảnh hoặc mặc cảm cố định.}

\vocabentry{Frustration}
{/frʌˈstreɪʃn/ (n)}
{The feeling of being upset or annoyed.}
{Lack of progress can often lead to feelings of intense frustration.}
{Đây là một danh từ dùng để chỉ sự thất vọng hoặc bực bội.}

\vocabentry{Function}
{/ˈfʌŋkʃn/ (n/v)}
{An activity or purpose natural to or intended for a person or thing.}
{One of the main functions of the brain is to process information from the senses.}
{Chức năng tâm lý.}

\vocabentry{Genetics}
{/dʒəˈnetɪks/ (n)}
{The study of heredity and the variation of inherited characteristics.}
{Genetics play a significant role in determining a person's temperament and intelligence.}
{Di truyền học ảnh hưởng đến tính cách.}

\vocabentry{Genius}
{/ˈdʒiːniəs/ (n)}
{Exceptional intellectual or creative power or other natural ability.}
{Albert Einstein was widely recognized as a scientific genius.}
{Đây là một danh từ dùng để chỉ thiên tài.}

\vocabentry{Gestalt}
{/ɡəˈʃtælt/ (n)}
{An organized whole that is perceived as more than the sum of its parts.}
{Gestalt psychology emphasizes that the human mind perceives patterns and wholes.}
{Đây là một danh từ dùng để chỉ thuyết Gestalt - tâm lý học hình thái.}

\vocabentry{Habit}
{/ˈæbɪt/ (n)}
{A settled or regular tendency or practice.}
{It takes about two months to form a new healthy habit like daily exercise.}
{Đây là một danh từ dùng để chỉ thói quen.}

\vocabentry{Heuristic}
{/hjuˈrɪstɪk/ (n)}
{A mental shortcut that allows people to solve problems and make judgments quickly.}
{Heuristics are useful for making quick decisions but can sometimes lead to errors in judgment.}
{Đây là một danh từ dùng để chỉ các lối tắt tư duy/quy tắc ngón tay cái.}

\vocabentry{Hierarchy}
{/ˈhaɪərɑːrki/ (n)}
{A system or organization in which people or groups are ranked one above the other according to status or authority.}
{Maslow's hierarchy of needs describes the different stages of human motivation.}
{Thứ bậc, như trong "Hierarchy of Needs".}

\vocabentry{Hostility}
{/hɑːˈstɪləti/ (n)}
{Unfriendly or hostile state, attitude, or action.}
{There was a sense of hidden hostility between the two rival teams.}
{Đây là một danh từ dùng để chỉ sự thù địch.}

\vocabentry{Humanism}
{/ˈhjuːmənɪzəm/ (n)}
{A rationalistic outlook or system of thought attaching prime importance to human rather than divine or supernatural matters.}
{Humanistic psychology emphasizes the importance of personal growth and self-actualization.}
{Đây là một danh từ dùng để chỉ chủ nghĩa nhân văn.}

\vocabentry{Hypnosis}
{/hɪpˈnoʊsɪs/ (n)}
{The induction of a state of consciousness in which a person apparently loses the power of voluntary action.}
{Hypnosis is sometimes used in therapy to help people overcome phobias or stop smoking.}
{Đây là một danh từ dùng để chỉ thôi miên.}

\vocabentry{Hysteria}
{/hɪˈstɪriə/ (n)}
{Exaggerated or uncontrollable emotion or excitement.}
{The news of the impending disaster caused widespread hysteria in the city.}
{Đây là một danh từ dùng để chỉ sự cuồng loạn hoặc kích động quá mức.}

\vocabentry{Identity}
{/aɪˈdentəti/ (n)}
{The fact of being who or what a person or thing is.}
{Adolescence is a key period for the development of a person's sense of identity.}
{Đây là một danh từ dùng để chỉ bản sắc hoặc danh tính cá nhân.}

\vocabentry{Ideology}
{/ˌaɪdiˈɑːlədʒi/ (n)}
{A system of ideas and ideals.}
{A person's political ideology can be influenced by their family and education.}
{Hệ tư tưởng.}

\vocabentry{Idiot}
{/ˈɪdiət/ (n)}
{A stupid person.}
{He felt like an idiot after making such a simple mistake.}
{Dùng trong lịch sử tâm lý học để chỉ mức độ thiểu năng, nay mang nghĩa xúc phạm.}

\vocabentry{Illusion}
{/ɪˈluːʒn/ (n)}
{A thing that is or is likely to be wrongly perceived or interpreted by the senses.}
{Optical illusions show how easily our brains can be tricked by visual information.}
{Đây là một danh từ dùng để chỉ ảo giác hoặc ảo tưởng.}

\vocabentry{Image}
{/ˈɪmɪdʒ/ (n)}
{A representation of the external form of a person or thing in art.}
{Having a negative self-image can contribute to eating disorders and depression.}
{Trong tâm lý là hình ảnh tự thân - "self-image".}

\vocabentry{Imagination}
{/ɪˌmædʒɪˈneɪʃn/ (n)}
{The faculty or action of forming new ideas, or images or concepts of external objects not present to the senses.}
{Children often have very vivid imaginations and enjoy playing make-believe games.}
{Đây là một danh từ dùng để chỉ trí tưởng tượng.}

\vocabentry{Imitation}
{/ˌɪmɪˈteɪʃn/ (n)}
{The action of using someone or something as a model.}
{Young children learn many skills through the imitation of their parents.}
{Sự bắt chước hành vi.}

\vocabentry{Impulse}
{/ˈɪmpʌls/ (n)}
{A sudden strong and unreflective urge or desire to act.}
{He had a sudden impulse to quit his job and travel the world.}
{Đây là một danh từ dùng để chỉ sự bốc đồng hoặc thôi thúc tức thời.}

\vocabentry{Individuality}
{/ˌɪndɪˌvɪdʒuˈæləti/ (n)}
{The quality or character of a particular person or thing that distinguishes them from others of the same kind.}
{Modern education should encourage students to express their own individuality.}
{Đây là một danh từ dùng để chỉ tính cá nhân hoặc cá tính.}

\vocabentry{Inequality}
{/ˌɪnɪˈkwɑːləti/ (n)}
{Difference in size, degree, circumstances, etc.; lack of equality.}
{Social inequality can lead to feelings of resentment and frustration among marginalized groups.}
{Sự bất bình đẳng ảnh hưởng đến tâm lý xã hội.}

\vocabentry{Insight}
{/ˈaɪsaɪt/ (n)}
{A deep understanding of a person or thing.}
{Therapy can provide people with valuable insight into their own behavior and emotions.}
{Sự thấu hiểu bản thân hoặc vấn đề.}

\vocabentry{Instinct}
{/ˈɪnstɪŋkt/ (n)}
{An innate, typically fixed pattern of behavior in animals in response to certain stimuli.}
{Birds have an amazing instinct for migration, traveling thousands of miles every year.}
{Đây là một danh từ dùng để chỉ bản năng.}

\vocabentry{Intelligence}
{/ɪnˈtelɪdʒəns/ (n)}
{The ability to acquire and apply knowledge and skills.}
{Intelligence is not just about doing well in school; it also involves problem-solving and social skills.}
{Đây là một danh từ dùng để chỉ trí thông minh.}

\vocabentry{Intention}
{/ɪnˈtenʃn/ (n)}
{A thing intended; an aim or plan.}
{He had every intention of finishing the project on time, but something unexpected happened.}
{Ý định hoặc mục đích hành động.}

\vocabentry{Interaction}
{/ˌɪntərˈækʃn/ (n)}
{Reciprocal action or influence.}
{Positive social interaction is essential for our mental and emotional well-being.}
{Sự tương tác xã hội.}

\vocabentry{Internalization}
{/ɪnˌtɜːrnələˈzeɪʃn/ (n)}
{The process of making attitudes or behavior part of one's own nature by learning or unconscious assimilation.}
{Children internalize the values and beliefs of their parents as they grow up.}
{Đây là một danh từ dùng để chỉ sự nội tâm hóa các giá trị.}

\vocabentry{Interpretation}
{/ɪnˌtɜːrprɪˈteɪʃn/ (n)}
{The action of explaining the meaning of something.}
{There can be many different interpretations of a single Rorschach inkblot.}
{Sự diễn giải hiện tượng tâm lý.}

\vocabentry{Introspection}
{/ˌɪntrəˈspekʃn/ (n)}
{The examination or observation of one's own mental and emotional processes.}
{Introspection can be a useful tool for personal growth and self-discovery.}
{Đây là một danh từ dùng để chỉ sự tự soi xét nội tâm.}

\vocabentry{Intuition}
{/ˌɪntuˈɪʃn/ (n)}
{The ability to understand something immediately, without the need for conscious reasoning.}
{Sometimes you just have to trust your intuition when making a difficult decision.}
{Đây là một danh từ dùng để chỉ trực giác.}

\vocabentry{Inversion}
{/ɪnˈvɜːrʒn/ (n)}
{The action of inverting something or the state of being inverted.}
{Freud used the term "inversion" to describe what we now call homosexuality.}
{Trong tâm lý học cổ điển dùng để chỉ sự đảo ngược xu hướng tình dục.}

\vocabentry{Isolation}
{/ˌaɪsəˈleɪʃn/ (n)}
{The process or fact of isolating or being isolated.}
{Long periods of isolation can have a negative impact on a person's mental health.}
{Sự cô lập xã hội.}

\vocabentry{Justice}
{/ˈdʒʌstɪs/ (n)}
{Just behavior or treatment.}
{A sense of justice is developed early in childhood as children learn to share and follow rules.}
{Cảm giác về sự công bằng.}

\vocabentry{Knowledge}
{/ˈnɑːlɪdʒ/ (n)}
{Facts, information, and skills acquired by a person through experience or education.}
{Having a basic knowledge of psychology can help you understand yourself and others better.}
{Tri thức về tâm lý.}

\vocabentry{Latency}
{/ˈleɪtənsi/ (n)}
{The state of existing but not yet being developed or manifest.}
{In Freud's theory, the latency period is a stage of childhood where sexual impulses are suppressed.}
{Giai đoạn tiềm ẩn/ủ bệnh.}

\vocabentry{Learning}
{/ˈlɜːrnɪŋ/ (n)}
{The acquisition of knowledge or skills through experience, study, or by being taught.}
{Observational learning occurs when we watch others and model our behavior on theirs.}
{Quá trình học tập và thay đổi hành vi.}

\vocabentry{Logic}
{/ˈlɑːdʒɪk/ (n)}
{Reasoning conducted or assessed according to strict principles of validity.}
{He used logic and reason to solve the complex puzzle.}
{Tư duy logic.}

\vocabentry{Maintenance}
{/ˈmeɪntənəns/ (n)}
{The process of maintaining or preserving someone or something.}
{Regular meditation is important for the maintenance of emotional balance.}
{Sự duy trì các thói quen tâm lý.}

\vocabentry{Maladjustment}
{/ˌmæləˈdʒʌstmənt/ (n)}
{Failure to cope with the demands of a normal social environment.}
{Maladjustment in school can lead to behavioral problems and academic failure.}
{Sự kém thích nghi xã hội.}

\vocabentry{Management}
{/ˈmænɪdʒmənt/ (n)}
{The process of dealing with or controlling things or people.}
{Stress management techniques can help people cope with the pressures of modern life.}
{Sự quản lý cảm xúc hoặc hành vi.}

\vocabentry{Masochistic}
{/ˌmæsəˈkɪstɪk/ (adj)}
{Deriving pleasure from one's own pain or humiliation.}
{Her masochistic tendencies led her to stay in an abusive relationship for years.}
{Đây là một tính từ dùng để chỉ tính tự ngược đãi.}

\vocabentry{Material}
{/məˈtɪriəl/ (n/adj)}
{The matter from which a thing is or can be made.}
{Excessive focus on material possessions can often lead to feelings of emptiness and dissatisfaction.}
{Vật chất ảnh hưởng đến tâm lý - "materialism".}

\vocabentry{Measure}
{/ˈmeʒər/ (n/v)}
{Ascertain the size, amount, or degree of (something) by using an instrument or device.}
{The IQ test is a common way to measure a person's intellectual potential.}
{Sự đo lường tâm lý.}

\vocabentry{Mechanism}
{/ˈmekənɪzəm/ (n)}
{A system of parts working together in a machine; a natural or established process by which something takes place.}
{The human brain has complex mechanisms for processing and storing information.}
{Cơ chế tâm lý.}

\vocabentry{Mediation}
{/ˌmiːdiˈeɪʃn/ (n)}
{Intervention in a dispute in order to resolve it.}
{Mediation can be a very effective way to resolve conflicts between family members.}
{Sự hòa giải xung đột.}

\vocabentry{Medicine}
{/ˈmedɪsn/ (n)}
{The science or practice of the diagnosis, treatment, and prevention of disease.}
{Modern medicine has developed many effective treatments for mental illnesses like depression.}
{Y học, bao gồm cả tâm thần học.}

\vocabentry{Memory}
{/ˈmeməri/ (n)}
{The faculty by which the mind stores and remembers information.}
{Short-term memory allows us to remember information for a few seconds or minutes.}
{Trí nhớ.}

\vocabentry{Mental}
{/ˈmentl/ (adj)}
{Relating to the mind.}
{Exercise is good for both your physical and mental health.}
{Đây là một tính từ dùng để chỉ những gì thuộc về tâm thần/tâm trí.}

\vocabentry{Metaphor}
{/ˈmetəfɔːr/ (n)}
{A figure of speech in which a word or phrase is applied to an object or action.}
{Using the metaphor of a "dark cloud" is a common way to describe feelings of depression.}
{Phép ẩn dụ mô tả tâm trạng.}

\vocabentry{Method}
{/ˈmeθəd/ (n)}
{A particular form of procedure for accomplishing or approaching something.}
{The scientific method is used to ensure that psychological research is reliable and valid.}
{Phương pháp nghiên cứu tâm lý.}

\vocabentry{Midlife}
{/ˌmɪdˈlaɪf/ (n)}
{The central period of a person's life, between around 45 and 60.}
{Many people experience a period of self-reflection and change during midlife.}
{Trung niên.}

\vocabentry{Mind}
{/maɪnd/ (n/v)}
{The element of a person that enables them to be aware of the world and their experiences.}
{Mindfulness involves paying attention to the present moment without judgment.}
{Tâm trí.}

\vocabentry{Minority}
{/maɪˈnɔːrəti/ (n)}
{The smaller number or part.}
{Minority groups often face discrimination and prejudice, which can impact their mental health.}
{Các nhóm thiểu số chịu áp lực tâm lý xã hội.}

\vocabentry{Mistake}
{/mɪˈsteɪk/ (n/v)}
{An action or judgment that is misguided or wrong.}
{It is important to learn from your mistakes and use them as an opportunity for growth.}
{Sai lầm trong nhận thức.}

\vocabentry{Model}
{/ˈmɑːdl/ (n/v)}
{A three-dimensional representation of a person or thing.}
{Parents should strive to be positive role models for their children.}
{Hình mẫu hành vi.}

\vocabentry{Mood}
{/muːd/ (n)}
{A temporary state of mind or feeling.}
{Her mood changed suddenly from happy to sad for no apparent reason.}
{Tâm trạng.}

\vocabentry{Moral}
{/ˈmɔːrəl/ (adj/n)}
{Concerned with the principles of right and wrong behavior.}
{Developing a strong moral compass is an important part of a child's upbringing.}
{Đạo đức cá nhân.}

\vocabentry{Motive}
{/ˈmoʊtɪv/ (n)}
{A reason for doing something.}
{The police are trying to determine the killer's motive for the crime.}
{Động cơ.}

\vocabentry{Mourning}
{/ˈmɔːrnɪŋ/ (n)}
{The expression of deep sorrow for someone who has died.}
{Mourning is a natural and necessary process after losing a loved one.}
{Sự đau buồn/tang chế.}

\vocabentry{Movement}
{/ˈmuːvmənt/ (n)}
{An act of changing physical location or position or of having this changed.}
{The positive psychology movement focuses on human strengths and well-being.}
{Phong trào tâm lý học.}

\vocabentry{Myth}
{/mɪθ/ (n)}
{A traditional story, especially one concerning the early history of a people.}
{Many ancient myths explore universal human themes and psychological conflicts.}
{Huyền thoại phản ánh tâm lý cộng đồng.}

\vocabentry{Narcissistic}
{/ˌnɑːrsɪˈsɪstɪk/ (adj)}
{Having an excessive interest in or admiration of oneself.}
{Narcissistic personality disorder is characterized by a lack of empathy and a constant need for admiration.}
{Tính ái kỷ.}

\vocabentry{Nature}
{/ˈneɪtʃər/ (n)}
{The basic or inherent features of something.}
{The "nature versus nurture" debate explores the relative influence of genetics and environment on behavior.}
{Bản chất/Tự nhiên.}

\vocabentry{Necessity}
{/nəˈsesəti/ (n)}
{The fact of being required or indispensable.}
{Sleep is a biological necessity for both physical and mental health.}
{Nhu cầu tất yếu của con người.}

\vocabentry{Need}
{/niːd/ (n/v)}
{Require (something) because it is essential or very important.}
{We all have a basic human need for social connection and belonging.}
{Nhu cầu tâm lý.}

\vocabentry{Negotiation}
{/nəˌɡoʊʃiˈeɪʃn/ (n)}
{Discussion aimed at reaching an agreement.}
{Good negotiation skills can help to resolve conflicts and build stronger relationships.}
{Sự đàm phán giải quyết xung đột.}

\vocabentry{Neurosis}
{/njʊˈroʊsɪs/ (n)}
{A relatively mild mental illness that is not caused by organic disease.}
{Anxiety neurosis can cause people to feel constant worry and fear.}
{Chứng loạn thần kinh chức năng.}

\vocabentry{Neutral}
{/ˈnuːtrəl/ (adj)}
{Not supporting or helping either side in a conflict, disagreement, etc.}
{A therapist should remain neutral and objective when working with clients.}
{Thái độ trung lập.}

\vocabentry{Norm}
{/nɔːrm/ (n)}
{Something that is usual, typical, or standard.}
{Breaking a social norm can sometimes lead to feelings of guilt or embarrassment.}
{Chuẩn mực xã hội.}

\vocabentry{Notion}
{/ˈnoʊʃn/ (n)}
{A conception of or belief about something.}
{The notion of "free will" is a central topic of debate in both philosophy and psychology.}
{Quan niệm tâm lý.}

\vocabentry{Nuance}
{/ˈnuːɑːns/ (n)}
{A subtle difference in or shade of meaning, expression, or sound.}
{A good psychologist is able to capture the subtle nuances of a client's emotions.}
{Sắc thái tâm lý tinh tế.}

\vocabentry{Nurture}
{/ˈnɜːrtʃər/ (n/v)}
{Care for and encourage the growth or development of.}
{A child's personality is shaped by both their genetics (nature) and their upbringing (nurture).}
{Sự nuôi dưỡng giáo dục.}

\vocabentry{Object}
{/ˈɑːbdʒekt/ (n)}
{A material thing that can be seen and touched.}
{For a small child, a favorite toy can become an important transitional object.}
{Trong tâm lý học là đối tượng gắn kết.}

\vocabentry{Objective}
{/əbˈdʒektɪv/ (adj/n)}
{(Of a person or their judgment) not influenced by personal feelings or opinions.}
{It is important to remain objective when analyzing scientific data.}
{Tính khách quan.}

\vocabentry{Observation}
{/ˌɑːbzərˈveɪʃn/ (n)}
{The action or process of observing something or someone carefully.}
{Naturalistic observation involves watching animals or people in their natural environment.}
{Sự quan sát hành vi.}

\vocabentry{Obsessive}
{/əbˈsesɪv/ (adj)}
{Of the nature of an obsession.}
{Obsessive-compulsive disorder (OCD) involves repetitive thoughts and behaviors.}
{Có tính chất ám ảnh.}

\vocabentry{Oedipus complex}
{/ˈedɪpəs ˈkɑːmpleks/ (n)}
{A Freudian theory of a child's unconscious sexual desire for the opposite-sex parent.}
{The Oedipus complex is one of the most famous and controversial theories in psychoanalysis.}
{Mặc cảm Oedipus.}

\vocabentry{Opinion}
{/əˈpɪnjən/ (n)}
{A view or judgment formed about something.}
{Everyone is entitled to their own opinion, even if you don't agree with it.}
{Ý kiến/Quan điểm cá nhân.}

\vocabentry{Oppression}
{/əˈpreʃn/ (n)}
{Prolonged cruel or unjust treatment or control.}
{Long-term social oppression can lead to feelings of hopelessness and despair.}
{Sự áp bức ảnh hưởng đến lòng tự trọng.}

\vocabentry{Option}
{/ˈɑːpʃn/ (n)}
{A thing that is or may be chosen.}
{When facing a difficult situation, it's important to consider all your options carefully.}
{Lựa chọn hành vi.}

\vocabentry{Oral}
{/ˈɔːrəl/ (adj)}
{Relating to the mouth.}
{Freud's "oral stage" is the first stage of psychosexual development in infants.}
{Giai đoạn môi miệng trong thuyết của Freud.}

\vocabentry{Order}
{/ˈɔːrdər/ (n/v)}
{The arrangement or disposition of people or things in relation to each other.}
{Some people have a strong need for order and predictability in their daily lives.}
{Trật tự hoặc rối loạn.}

\vocabentry{Organization}
{/ˌɔːrɡənəˈzeɪʃn/ (n)}
{An organized body of people with a particular purpose.}
{The internal organization of a person's personality is complex and unique.}
{Sự tổ chức tâm trí/nhân cách.}

\vocabentry{Orientation}
{/ˌɔːriənˈteɪʃn/ (n)}
{The action of orienting someone or something relative to specified positions.}
{Sexual orientation is an enduring pattern of emotional and romantic attraction.}
{Sự định hướng - ví dụ: sexual orientation.}

\vocabentry{Outcome}
{/ˈaʊtkʌm/ (n)}
{The way a thing turns out; a consequence.}
{The outcome of the experiment confirmed the researchers' hypothesis.}
{Kết quả của một quá trình tâm lý.}

\vocabentry{Outline}
{/ˈaʊtlaɪn/ (n/v)}
{A general description or plan giving the essential features of something.}
{The psychologist provided a brief outline of the patient's treatment plan.}
{Dàn ý/Phác thảo tính cách.}

\vocabentry{Overcome}
{/ˌoʊvərˈkʌm/ (v)}
{Succeed in dealing with a problem or difficulty.}
{With hard work and support, she was able to overcome her fear of public speaking.}
{Vượt qua nỗi sợ/khó khăn.}

\vocabentry{Panic}
{/ˈpænɪk/ (n/v)}
{Sudden uncontrollable fear or anxiety.}
{Panic attacks can cause physical symptoms like heart palpitations and shortness of breath.}
{Sự hoảng loạn.}

\vocabentry{Paradigm}
{/ˈpærədaɪm/ (n)}
{A typical example or pattern of something; a model.}
{The cognitive revolution created a new paradigm for psychological research.}
{Mô hình/Khuôn mẫu tư duy.}

\vocabentry{Paranoid}
{/ˈpærənɔɪd/ (adj)}
{Characterized by delusions of persecution or unwarranted jealousy.}
{He became increasingly paranoid and started to believe that his coworkers were out to get him.}
{Có tính hoang tưởng.}

\vocabentry{Participation}
{/pɑːrˌtɪsɪˈpeɪʃn/ (n)}
{The action of taking part in something.}
{Active participation in community activities can improve a person's sense of belonging.}
{Sự tham gia xã hội.}

\vocabentry{Passion}
{/ˈpæʃn/ (n)}
{Strong and barely controllable emotion.}
{He has a great passion for music and spends most of his free time playing the guitar.}
{Niềm đam mê/Cảm xúc nồng nhiệt.}

\vocabentry{Passive}
{/ˈpæsɪv/ (adj)}
{Accepting or allowing what happens or what others do, without active response.}
{A passive communication style involves failing to express your own needs and feelings.}
{Tính thụ động.}

\vocabentry{Pathology}
{/pəˈθɑːlədʒi/ (n)}
{The science of the causes and effects of diseases.}
{Psychopathology is the study of mental disorders and their causes.}
{Bệnh lý học tâm thần.}

\vocabentry{Pattern}
{/ˈpætərn/ (n/v)}
{A repeated decorative design; a regular and intelligible form or sequence.}
{Recognizing patterns of negative thinking is the first step towards changing them.}
{Khuôn mẫu hành vi.}

\vocabentry{Pedagogy}
{/ˈpedəɡɑːdʒi/ (n)}
{The method and practice of teaching.}
{Effective pedagogy takes into account the different learning styles of students.}
{Tâm lý giáo dục.}

\vocabentry{Peer}
{/pɪr/ (n/v)}
{A person of the same age, status, or ability as another specified person.}
{Peer pressure can have a strong influence on the behavior of adolescents.}
{Bạn đồng trang lứa.}

\vocabentry{Penalty}
{/ˈpenəlti/ (n)}
{A punishment imposed for breaking a law, rule, or contract.}
{In behavioral psychology, a penalty is used to decrease the frequency of an unwanted behavior.}
{Hình phạt trong điều kiện hóa hành vi.}

\vocabentry{Perceive}
{/pərˈsiːv/ (v)}
{Become aware or conscious of (something); come to realize or understand.}
{People often perceive the same event in very different ways.}
{Nhận thức được.}

\vocabentry{Performance}
{/pərˈfɔːrməns/ (n)}
{An act of staging or presenting a play, concert, or other form of entertainment.}
{Anxiety can have a negative impact on a person's performance in high-pressure situations.}
{Hiệu suất thực hiện công việc/hành vi.}

\vocabentry{Period}
{/ˈpɪriəd/ (n)}
{A length or portion of time.}
{Adolescence is a critical period for the development of social and emotional skills.}
{Giai đoạn phát triển.}

\vocabentry{Personality}
{/ˌpɜːrsəˈnæləti/ (n)}
{The combination of characteristics or qualities that form an individual's distinctive character.}
{Your personality is shaped by a complex interplay of genetics and environment.}
{Đây là một danh từ dùng để chỉ nhân cách hoặc tính cách cá nhân.}

